\chapter{NLP Architectures}
\label{chap:nlp_architectures}

\begin{tldr}
This chapter covers the three major Transformer paradigms for NLP: encoder-only (BERT family), decoder-only (GPT family), and encoder-decoder (T5/BART). Understanding when to use each architecture—and the design decisions within each family—is a common staff-level interview topic. We also cover tokenization methods (BPE, WordPiece, SentencePiece), modern LLM innovations, and practical guidance for architecture selection.
\end{tldr}

\section{The Three Transformer Paradigms}
\label{sec:transformer_paradigms}

\subsection{Architectural Overview}

\begin{figure}[H]
\centering
\begin{tikzpicture}[
    node distance=0.8cm,
    box/.style={rectangle, draw, minimum width=2.5cm, minimum height=0.7cm, rounded corners},
    encoder/.style={box, fill=lightblue},
    decoder/.style={box, fill=lightgreen},
    cross/.style={box, fill=orange!30}
]
    % Encoder-only
    \begin{scope}[shift={(0,0)}]
        \node[encoder] (e1) {Encoder};
        \node[encoder, above=0.3cm of e1] (e2) {Encoder};
        \node[above=0.3cm of e2] (edots) {$\vdots$};
        \node[encoder, above=0.3cm of edots] (eN) {Encoder};
        \node[below=0.5cm of e1] {Input};
        \node[above=0.5cm of eN] {[CLS] / Tokens};
        \node[below=1.2cm of e1, font=\bfseries] {Encoder-Only};
        \node[below=1.7cm of e1, font=\small] {(BERT)};
        
        % Bidirectional arrows
        \draw[<->, thick, blue] ($(e1.west)+(-0.3,0)$) -- ($(eN.west)+(-0.3,0)$);
        \node[left=0.5cm of e2, font=\tiny, blue] {Bidirectional};
    \end{scope}
    
    % Decoder-only
    \begin{scope}[shift={(5,0)}]
        \node[decoder] (d1) {Decoder};
        \node[decoder, above=0.3cm of d1] (d2) {Decoder};
        \node[above=0.3cm of d2] (ddots) {$\vdots$};
        \node[decoder, above=0.3cm of ddots] (dN) {Decoder};
        \node[below=0.5cm of d1] {Input};
        \node[above=0.5cm of dN] {Next Token};
        \node[below=1.2cm of d1, font=\bfseries] {Decoder-Only};
        \node[below=1.7cm of d1, font=\small] {(GPT)};
        
        % Causal arrow
        \draw[->, thick, green!70!black] ($(d1.east)+(0.3,0)$) -- ($(dN.east)+(0.3,0)$);
        \node[right=0.5cm of d2, font=\tiny, green!70!black] {Causal};
    \end{scope}
    
    % Encoder-Decoder
    \begin{scope}[shift={(10,0)}]
        \node[encoder] (ed_e1) {Encoder};
        \node[encoder, above=0.3cm of ed_e1] (ed_eN) {Encoder};
        
        \node[decoder, right=1cm of ed_e1] (ed_d1) {Decoder};
        \node[cross, above=0.3cm of ed_d1] (ed_cross) {Cross-Attn};
        \node[decoder, above=0.3cm of ed_cross] (ed_dN) {Decoder};
        
        \node[below=0.5cm of ed_e1] {Source};
        \node[below=0.5cm of ed_d1] {Target};
        \node[above=0.5cm of ed_dN] {Output};
        \node[below=1.2cm of $(ed_e1)!0.5!(ed_d1)$, font=\bfseries] {Encoder-Decoder};
        \node[below=1.7cm of $(ed_e1)!0.5!(ed_d1)$, font=\small] {(T5, BART)};
        
        % Cross attention arrow
        \draw[->, thick, orange] (ed_eN.east) -- (ed_cross.west);
    \end{scope}
\end{tikzpicture}
\caption{The three Transformer paradigms for NLP.}
\label{fig:transformer_paradigms}
\end{figure}

\subsection{Key Differences}

\begin{table}[H]
\centering
\caption{Comparison of Transformer paradigms}
\begin{tabularx}{\textwidth}{lXXX}
\toprule
\textbf{Aspect} & \textbf{Encoder-Only} & \textbf{Decoder-Only} & \textbf{Encoder-Decoder} \\
\midrule
Attention & Bidirectional & Causal (left-to-right) & Bidirectional enc, Causal dec \\
Context & Full sequence & Past tokens only & Full source, past target \\
Pretraining & Masked LM (MLM) & Causal LM & Span corruption / denoising \\
Output & Representations & Next token prob & Next token prob \\
Best for & Understanding & Generation & Seq-to-seq \\
Examples & BERT, RoBERTa & GPT, LLaMA & T5, BART \\
\bottomrule
\end{tabularx}
\end{table}

%==============================================================================
\section{Encoder-Only Models (BERT Family)}
\label{sec:encoder_only}
%==============================================================================

\subsection{BERT: Bidirectional Encoder Representations from Transformers}

\subsubsection{Architecture}

\begin{itemize}
    \item Stack of Transformer encoder blocks
    \item Bidirectional self-attention (no causal mask)
    \item Special tokens: [CLS] for classification, [SEP] for segment separation
    \item Segment embeddings for sentence pairs
\end{itemize}

\subsubsection{Pre-training Objectives}

\textbf{Masked Language Modeling (MLM)}:
\begin{itemize}
    \item Randomly mask 15\% of tokens
    \item Of masked positions:
    \begin{itemize}
        \item 80\%: Replace with [MASK]
        \item 10\%: Replace with random token
        \item 10\%: Keep original
    \end{itemize}
    \item Predict original token
\end{itemize}

\begin{equation}
\mathcal{L}_{MLM} = -\E\left[\sum_{i \in \mathcal{M}} \log p(x_i | \tilde{\vx})\right]
\end{equation}

where $\mathcal{M}$ is the set of masked positions and $\tilde{\vx}$ is the corrupted input.

\textbf{Why the 80/10/10 split?}
\begin{itemize}
    \item [MASK] never appears at fine-tuning time
    \item Training only on [MASK] creates mismatch
    \item Random replacement forces model to maintain representations for all tokens
    \item Original token prevents trivial copying
\end{itemize}

\textbf{Next Sentence Prediction (NSP)}:
\begin{itemize}
    \item Given [CLS] Sentence A [SEP] Sentence B [SEP]
    \item Binary classification: Is B the actual next sentence?
    \item 50\% positive (actual next), 50\% negative (random)
\end{itemize}

Note: NSP was later found to be less important; RoBERTa removes it.

\subsubsection{Fine-tuning Patterns}

\begin{figure}[H]
\centering
\begin{tikzpicture}[
    node distance=0.5cm,
    box/.style={rectangle, draw, minimum width=1.5cm, minimum height=0.5cm, rounded corners, font=\small},
    bert/.style={box, fill=lightblue, minimum width=4cm}
]
    % Classification
    \begin{scope}[shift={(0,0)}]
        \node[bert] (bert1) {BERT};
        \node[box, fill=lightgreen, above=0.3cm of bert1] (cls1) {Classifier};
        \node[below=0.3cm of bert1, font=\small] {[CLS] sent [SEP]};
        \node[above=0.3cm of cls1, font=\small] {Label};
        \draw[->] (bert1) -- (cls1);
        \node[below=1cm of bert1, font=\bfseries\small] {Classification};
    \end{scope}
    
    % NER
    \begin{scope}[shift={(5,0)}]
        \node[bert] (bert2) {BERT};
        \node[box, fill=lightgreen, above=0.3cm of bert2, minimum width=4cm] (ner) {Token Classifier};
        \node[below=0.3cm of bert2, font=\small] {$t_1$ $t_2$ ... $t_n$};
        \node[above=0.3cm of ner, font=\small] {$l_1$ $l_2$ ... $l_n$};
        \draw[->] (bert2) -- (ner);
        \node[below=1cm of bert2, font=\bfseries\small] {Token Classification};
    \end{scope}
    
    % QA
    \begin{scope}[shift={(10,0)}]
        \node[bert] (bert3) {BERT};
        \node[box, fill=lightgreen, above=0.3cm of bert3, minimum width=4cm] (qa) {Start/End Classifier};
        \node[below=0.3cm of bert3, font=\small] {[CLS] Q [SEP] Context};
        \node[above=0.3cm of qa, font=\small] {Start, End positions};
        \draw[->] (bert3) -- (qa);
        \node[below=1cm of bert3, font=\bfseries\small] {Extractive QA};
    \end{scope}
\end{tikzpicture}
\caption{Common BERT fine-tuning patterns.}
\label{fig:bert_finetuning}
\end{figure}

\subsection{BERT Variants}

\begin{table}[H]
\centering
\caption{BERT family comparison}
\begin{tabularx}{\textwidth}{lXXX}
\toprule
\textbf{Model} & \textbf{Key Changes} & \textbf{Improvement} & \textbf{Parameters} \\
\midrule
BERT-base & Original & Baseline & 110M \\
BERT-large & Larger & +2-3\% on tasks & 340M \\
RoBERTa & No NSP, more data, dynamic masking & +2-5\% & 125M / 355M \\
ALBERT & Factorized embeddings, cross-layer sharing & 18× fewer params & 12M - 235M \\
DeBERTa & Disentangled attention, enhanced mask decoder & SOTA on many benchmarks & 100M - 1.5B \\
ELECTRA & Replaced token detection (more efficient) & 4× sample efficient & 14M - 335M \\
\bottomrule
\end{tabularx}
\end{table}

\subsubsection{RoBERTa Improvements}

\begin{enumerate}
    \item \textbf{Remove NSP}: NSP task didn't help, possibly hurt
    \item \textbf{Dynamic masking}: Different mask each epoch (vs. static)
    \item \textbf{Larger batches}: 8K batch size
    \item \textbf{More data}: 160GB text (vs. BERT's 16GB)
    \item \textbf{Longer training}: 500K steps (vs. 1M, but larger batches)
\end{enumerate}

\subsubsection{ELECTRA: Efficient Pre-training}

\textbf{Key insight}: MLM only provides signal for 15\% of tokens.

\textbf{ELECTRA approach}:
\begin{enumerate}
    \item Train small generator to fill in [MASK] tokens
    \item Train discriminator to detect which tokens were replaced
    \item All tokens provide training signal
\end{enumerate}

\begin{equation}
\mathcal{L}_{ELECTRA} = \sum_{i=1}^{n} -\mathbb{1}[x_i^{corrupt} \neq x_i^{original}] \log D(x_i) - \mathbb{1}[x_i^{corrupt} = x_i^{original}] \log(1 - D(x_i))
\end{equation}

\textbf{Result}: Matches BERT performance with 1/4 the compute.

\subsubsection{DeBERTa: Disentangled Attention}

\textbf{Key innovation}: Separate content and position representations.

Standard attention:
\begin{equation}
\text{Attention} = \text{softmax}((\mathbf{x}_i + \mathbf{p}_i)\mathbf{W}^Q \cdot ((\mathbf{x}_j + \mathbf{p}_j)\mathbf{W}^K)^T)
\end{equation}

Disentangled:
\begin{equation}
\text{Attention} = \text{softmax}(\mathbf{x}_i\mathbf{W}^Q \cdot (\mathbf{x}_j\mathbf{W}^K)^T + \mathbf{x}_i\mathbf{W}^Q \cdot (\mathbf{p}_{i-j}\mathbf{W}^K)^T + \mathbf{p}_{i-j}\mathbf{W}^Q \cdot (\mathbf{x}_j\mathbf{W}^K)^T)
\end{equation}

Content-to-content + Content-to-position + Position-to-content

\subsection{When to Use Encoder-Only}

\begin{itemize}
    \item \textbf{Text classification}: Sentiment, topic, intent
    \item \textbf{Token classification}: NER, POS tagging
    \item \textbf{Extractive QA}: Find answer span in context
    \item \textbf{Semantic similarity}: Sentence embeddings, retrieval
    \item \textbf{NOT for}: Open-ended generation
\end{itemize}

%==============================================================================
\section{Decoder-Only Models (GPT Family)}
\label{sec:decoder_only}
%==============================================================================

\subsection{Architecture}

\begin{itemize}
    \item Stack of Transformer decoder blocks (no cross-attention)
    \item Causal self-attention (each position only attends to past)
    \item No [CLS] token; use last token or pooling for classification
    \item Autoregressive generation
\end{itemize}

\subsection{Causal Language Modeling}

\begin{mathresult}{Causal LM Objective}
\begin{equation}
\mathcal{L}_{CLM} = -\sum_{i=1}^{n} \log p(x_i | x_1, \ldots, x_{i-1}; \theta)
\end{equation}
\end{mathresult}

\textbf{Key property}: Every token provides a training signal (vs. 15\% for MLM).

\textbf{Inference}: Generate token by token, feeding each output back as input.

\subsection{Model Evolution}

\begin{table}[H]
\centering
\caption{GPT family evolution}
\begin{tabularx}{\textwidth}{lrrlX}
\toprule
\textbf{Model} & \textbf{Params} & \textbf{Context} & \textbf{Training Data} & \textbf{Key Innovation} \\
\midrule
GPT & 117M & 512 & BookCorpus & First large-scale CLM \\
GPT-2 & 1.5B & 1024 & WebText (40GB) & Zero-shot capabilities \\
GPT-3 & 175B & 2048 & 570GB & In-context learning \\
GPT-4 & ~1.8T? & 8K/32K & Unknown & Multimodal, RLHF \\
\bottomrule
\end{tabularx}
\end{table}

\subsection{Open Source LLMs}

\begin{table}[H]
\centering
\caption{Major open-source LLMs}
\begin{tabularx}{\textwidth}{lrrlX}
\toprule
\textbf{Model} & \textbf{Params} & \textbf{Context} & \textbf{Organization} & \textbf{Key Features} \\
\midrule
LLaMA & 7-65B & 2K & Meta & Efficient, open weights \\
LLaMA 2 & 7-70B & 4K & Meta & GQA, longer context \\
Mistral & 7B & 8K & Mistral AI & Sliding window attention \\
Mixtral & 8×7B & 32K & Mistral AI & MoE, 12B active \\
Falcon & 7-180B & 2K & TII & Multi-query attention \\
\bottomrule
\end{tabularx}
\end{table}

\subsection{Modern LLM Design Choices}

\subsubsection{LLaMA Architecture Details}

\begin{enumerate}
    \item \textbf{Pre-normalization}: RMSNorm before attention and FFN
    \item \textbf{SwiGLU activation}: $\text{FFN}(x) = (\text{Swish}(xW_1) \odot xV)W_2$
    \item \textbf{RoPE}: Rotary positional embeddings
    \item \textbf{No bias}: Remove bias terms throughout
    \item \textbf{GQA} (LLaMA 2): Grouped-query attention for efficiency
\end{enumerate}

\subsubsection{Mistral Innovations}

\begin{enumerate}
    \item \textbf{Sliding window attention}: Each layer attends to window of 4096 tokens
    \item \textbf{Rolling buffer cache}: Fixed memory regardless of sequence length
    \item \textbf{Pre-fill and chunking}: Efficient handling of long prompts
\end{enumerate}

\subsection{In-Context Learning}

\textbf{In-Context Learning.} The ability to perform tasks by conditioning on examples in the prompt, without gradient updates.

\textbf{Example (few-shot sentiment)}:
\begin{verbatim}
Review: "Great movie!" Sentiment: Positive
Review: "Terrible waste of time" Sentiment: Negative
Review: "Absolutely loved it!" Sentiment: [MODEL GENERATES]
\end{verbatim}

\textbf{Emergent at scale}: ICL doesn't appear in small models; emerges around 10B+ parameters.

\subsection{Generation Strategies}

\subsubsection{Greedy Decoding}
\begin{equation}
x_t = \argmax_x p(x | x_{<t})
\end{equation}
Simple but repetitive, lacks diversity.

\subsubsection{Beam Search}
Maintain $k$ best partial sequences.
\begin{itemize}
    \item Better for constrained generation (translation)
    \item Still tends toward repetition for open-ended generation
\end{itemize}

\subsubsection{Temperature Sampling}
\begin{equation}
p(x | x_{<t}) = \frac{\exp(z_x / T)}{\sum_{x'} \exp(z_{x'} / T)}
\end{equation}

\begin{itemize}
    \item $T < 1$: Sharper, more deterministic
    \item $T > 1$: Flatter, more random
    \item $T = 1$: Standard softmax
\end{itemize}

\subsubsection{Top-k Sampling}
Sample only from top-$k$ most likely tokens.

\subsubsection{Nucleus (Top-p) Sampling}
Sample from smallest set of tokens whose cumulative probability $\geq p$.

\begin{equation}
V_p = \argmin_V \left\{\sum_{x \in V} p(x | x_{<t}) \geq p\right\}
\end{equation}

\textbf{Advantage}: Adapts to probability distribution shape (more tokens when uncertain).

\subsubsection{Recommended Settings}

\begin{table}[H]
\centering
\caption{Generation settings by use case}
\begin{tabularx}{\textwidth}{lXX}
\toprule
\textbf{Use Case} & \textbf{Settings} & \textbf{Rationale} \\
\midrule
Factual QA & $T=0$ or greedy & Deterministic, accurate \\
Creative writing & $T=0.7-1.0$, top-p=0.9 & Diverse but coherent \\
Code generation & $T=0.2-0.5$, top-p=0.95 & Low diversity, high precision \\
Brainstorming & $T=1.0+$, top-p=0.95 & Maximum diversity \\
\bottomrule
\end{tabularx}
\end{table}

\subsection{When to Use Decoder-Only}

\begin{itemize}
    \item \textbf{Text generation}: Stories, articles, completions
    \item \textbf{Conversational AI}: Chatbots, assistants
    \item \textbf{Code generation}: Copilot-style applications
    \item \textbf{Few-shot learning}: When you can't fine-tune
    \item \textbf{General-purpose}: When one model must do many tasks
\end{itemize}

%==============================================================================
\section{Encoder-Decoder Models (T5 Family)}
\label{sec:encoder_decoder}
%==============================================================================

\subsection{Architecture}

\begin{itemize}
    \item Separate encoder and decoder stacks
    \item Encoder: Bidirectional self-attention
    \item Decoder: Causal self-attention + cross-attention to encoder
    \item Both encoder and decoder outputs used
\end{itemize}

\subsection{T5: Text-to-Text Transfer Transformer}

\subsubsection{Key Innovation: Unified Text-to-Text Format}

All tasks framed as text-to-text:
\begin{itemize}
    \item \textbf{Translation}: ``translate English to German: [text]'' → German text
    \item \textbf{Classification}: ``sentiment: [text]'' → ``positive'' / ``negative''
    \item \textbf{Summarization}: ``summarize: [text]'' → summary
    \item \textbf{QA}: ``question: [q] context: [c]'' → answer
\end{itemize}

\subsubsection{Pre-training: Span Corruption}

\begin{enumerate}
    \item Randomly select spans (average 3 tokens each)
    \item Replace each span with sentinel token ([MASK1], [MASK2], ...)
    \item Decoder outputs: sentinel + original span tokens
\end{enumerate}

\begin{example}
Input: ``The [MASK1] brown fox [MASK2] the lazy dog''\\
Target: ``[MASK1] quick [MASK2] jumped over''
\end{example}

\subsubsection{T5 Variants}

\begin{table}[H]
\centering
\caption{T5 model sizes}
\begin{tabular}{lrrr}
\toprule
\textbf{Model} & \textbf{Params} & \textbf{Layers} & \textbf{Dim} \\
\midrule
T5-Small & 60M & 6 enc + 6 dec & 512 \\
T5-Base & 220M & 12 + 12 & 768 \\
T5-Large & 770M & 24 + 24 & 1024 \\
T5-3B & 3B & 24 + 24 & 1024 \\
T5-11B & 11B & 24 + 24 & 1024 \\
\bottomrule
\end{tabular}
\end{table}

\subsection{BART: Denoising Autoencoder}

\subsubsection{Pre-training: Multiple Corruption Schemes}

\begin{enumerate}
    \item \textbf{Token masking}: Replace with [MASK]
    \item \textbf{Token deletion}: Remove tokens (no placeholder)
    \item \textbf{Text infilling}: Replace span with single [MASK]
    \item \textbf{Sentence permutation}: Shuffle sentences
    \item \textbf{Document rotation}: Rotate to start from random token
\end{enumerate}

\textbf{Best configuration}: Text infilling with 30\% span masking, Poisson span lengths ($\lambda = 3$).

\subsection{When to Use Encoder-Decoder}

\begin{itemize}
    \item \textbf{Translation}: Explicit alignment between source and target
    \item \textbf{Summarization}: Condense long source to short target
    \item \textbf{Structured generation}: When output structure differs from input
    \item \textbf{Conditional generation}: Strong conditioning on input needed
\end{itemize}

%==============================================================================
\section{Architecture Selection Guide}
\label{sec:nlp_arch_selection}
%==============================================================================

\begin{figure}[H]
\centering
\begin{tikzpicture}[
    node distance=1cm,
    decision/.style={diamond, draw, aspect=2.5, inner sep=1pt, fill=lightyellow, font=\small},
    result/.style={rectangle, draw, rounded corners, fill=lightgreen, minimum width=2cm, font=\small},
    arrow/.style={->, thick}
]
    % Start
    \node[decision] (gen) {Need generation?};
    
    % No generation
    \node[decision, below left=1.5cm and 1cm of gen] (und) {Task type?};
    \node[result, below left=1cm and 0cm of und] (bert) {BERT/RoBERTa};
    \node[result, below right=1cm and 0cm of und] (deberta) {DeBERTa};
    
    % Generation
    \node[decision, below right=1.5cm and 1cm of gen] (cond) {Conditional?};
    \node[result, below left=1cm and 0cm of cond] (t5) {T5/BART};
    \node[result, below right=1cm and 0cm of cond] (gpt) {GPT/LLaMA};
    
    % Edges
    \draw[arrow] (gen) -- node[above left, font=\small] {No} (und);
    \draw[arrow] (gen) -- node[above right, font=\small] {Yes} (cond);
    \draw[arrow] (und) -- node[left, font=\small] {Classification} (bert);
    \draw[arrow] (und) -- node[right, font=\small] {SOTA needed} (deberta);
    \draw[arrow] (cond) -- node[left, font=\small] {Yes} (t5);
    \draw[arrow] (cond) -- node[right, font=\small] {No/Few-shot} (gpt);
\end{tikzpicture}
\caption{NLP architecture selection decision tree.}
\label{fig:nlp_decision_tree}
\end{figure}

\begin{table}[H]
\centering
\caption{Architecture selection by task}
\begin{tabularx}{\textwidth}{lXX}
\toprule
\textbf{Task} & \textbf{Recommended} & \textbf{Alternative} \\
\midrule
Sentiment classification & BERT, RoBERTa & Fine-tuned LLM \\
Named Entity Recognition & BERT + CRF & Token classifier \\
Extractive QA & BERT, DeBERTa & - \\
Generative QA & T5, GPT & - \\
Machine translation & T5, mBART & - \\
Summarization & BART, T5 & GPT (zero-shot) \\
Dialogue / Chat & GPT, LLaMA & - \\
Code generation & CodeLLaMA, StarCoder & GPT-4 \\
Semantic similarity & Sentence-BERT & E5, BGE \\
Zero-shot classification & GPT, LLaMA & BART (NLI-based) \\
\bottomrule
\end{tabularx}
\end{table}

\begin{interviewq}
\textbf{Q: You need to build a system for customer support ticket classification with 50 categories. Which architecture would you choose and why?}

\textbf{A:} I would choose \textbf{BERT or RoBERTa} for several reasons:

\begin{enumerate}
    \item \textbf{Task type}: Classification is an understanding task—bidirectional context is valuable.
    
    \item \textbf{Efficiency}: Encoder-only models are faster for classification (single forward pass vs. generation).
    
    \item \textbf{Fine-tuning}: With 50 categories, we need task-specific fine-tuning. Encoder models are designed for this.
    
    \item \textbf{Latency}: Production classification needs fast inference—encoders are faster than decoders.
\end{enumerate}

\textbf{If I couldn't fine-tune} (e.g., no labeled data): I would use GPT/LLaMA with few-shot prompting or use a zero-shot approach with an NLI model (``This ticket is about [category]'' → entailment score).

\textbf{Architecture details}: RoBERTa-base → [CLS] embedding → Dropout → Linear(768, 50) → Softmax.
\end{interviewq}

%==============================================================================
\section{Tokenization}
\label{sec:tokenization}
%==============================================================================

\subsection{Why Tokenization Matters}

\begin{itemize}
    \item \textbf{Vocabulary size}: Affects embedding table size and rare word handling
    \item \textbf{Sequence length}: Character-level is longer, word-level shorter
    \item \textbf{OOV handling}: Subword methods handle unseen words gracefully
    \item \textbf{Multilingual}: Different scripts need different treatment
\end{itemize}

\subsection{Tokenization Methods}

\subsubsection{Byte Pair Encoding (BPE)}

\begin{algorithm}[H]
\caption{BPE Training}
\begin{algorithmic}[1]
\REQUIRE Corpus, vocabulary size $V$
\STATE Initialize vocabulary with all characters
\WHILE{$|$vocabulary$| < V$}
    \STATE Count frequency of all adjacent pairs
    \STATE Merge most frequent pair into new token
    \STATE Add new token to vocabulary
\ENDWHILE
\RETURN vocabulary, merge rules
\end{algorithmic}
\end{algorithm}

\textbf{Example}:
\begin{enumerate}
    \item Start: l, o, w, e, r, \_
    \item Most frequent pair: (e, r) → Merge to ``er''
    \item Next: (l, o) → ``lo''
    \item Continue until vocabulary size reached
\end{enumerate}

\textbf{Used in}: GPT-2, GPT-3, RoBERTa

\subsubsection{WordPiece}

Similar to BPE but uses likelihood-based merging:
\begin{equation}
\text{score}(x, y) = \frac{\text{freq}(xy)}{\text{freq}(x) \times \text{freq}(y)}
\end{equation}

Merge pair that maximizes likelihood of training data.

\textbf{Used in}: BERT, DistilBERT

\textbf{Notation}: Continuations marked with \#\# (e.g., ``playing'' → ``play'', ``\#\#ing'')

\subsubsection{Unigram}

\begin{enumerate}
    \item Start with large vocabulary
    \item Iteratively remove tokens that least affect likelihood
    \item Keep tokens that maximize: $\sum_{\text{corpus}} \log p(\text{tokenization})$
\end{enumerate}

\textbf{Advantage}: Can output multiple tokenizations with probabilities.

\textbf{Used in}: T5, ALBERT

\subsubsection{SentencePiece}

\begin{itemize}
    \item Language-agnostic (no pre-tokenization)
    \item Treats input as raw byte stream
    \item Uses BPE or Unigram internally
    \item Handles any script (CJK, emoji, etc.)
\end{itemize}

\textbf{Used in}: LLaMA, T5, mBERT

\subsection{Comparison}

\begin{table}[H]
\centering
\caption{Tokenization method comparison}
\begin{tabularx}{\textwidth}{lXXX}
\toprule
\textbf{Method} & \textbf{Pros} & \textbf{Cons} & \textbf{Used By} \\
\midrule
BPE & Efficient, deterministic & Greedy & GPT, RoBERTa \\
WordPiece & Likelihood-based & More complex & BERT \\
Unigram & Probabilistic, multiple tokenizations & Slower training & T5, ALBERT \\
SentencePiece & Language-agnostic, handles any script & — & LLaMA, T5 \\
\bottomrule
\end{tabularx}
\end{table}

\subsection{Practical Considerations}

\begin{itemize}
    \item \textbf{Vocabulary size}: 30K-50K typical; larger for multilingual
    \item \textbf{Byte-level BPE}: GPT-2 uses byte-level (256 base tokens) for complete coverage
    \item \textbf{Special tokens}: [PAD], [UNK], [CLS], [SEP], [MASK], etc.
    \item \textbf{Normalization}: Lowercasing, unicode normalization, etc.
\end{itemize}

%==============================================================================
\section{Instruction Tuning and RLHF}
\label{sec:instruction_tuning}
%==============================================================================

\subsection{The Alignment Problem}

Base language models are trained to predict next tokens, not to be helpful. They may:
\begin{itemize}
    \item Continue harmful text
    \item Refuse benign requests
    \item Not follow instructions
    \item Hallucinate confidently
\end{itemize}

\subsection{Instruction Tuning}

Fine-tune on instruction-following examples:
\begin{verbatim}
Instruction: Summarize the following article in one sentence.
Input: [article text]
Output: [one-sentence summary]
\end{verbatim}

\textbf{Key datasets}: FLAN, Natural Instructions, Alpaca, Dolly

\textbf{Effect}: Model learns to follow diverse instructions without task-specific fine-tuning.

\subsection{RLHF: Reinforcement Learning from Human Feedback}

\begin{figure}[H]
\centering
\begin{tikzpicture}[
    node distance=1cm,
    box/.style={rectangle, draw, minimum width=2.5cm, minimum height=0.8cm, rounded corners},
    stage/.style={box, fill=lightblue}
]
    % Stage 1
    \node[stage] (sft) {SFT Model};
    \node[below=0.3cm of sft, font=\small] {Supervised Fine-Tuning};
    
    % Stage 2
    \node[stage, right=2cm of sft] (rm) {Reward Model};
    \node[below=0.3cm of rm, font=\small] {Train on preferences};
    
    % Stage 3
    \node[stage, right=2cm of rm] (ppo) {RLHF Model};
    \node[below=0.3cm of ppo, font=\small] {PPO optimization};
    
    % Edges
    \draw[->] (sft) -- (rm) node[midway, above, font=\small] {Generate};
    \draw[->] (rm) -- (ppo) node[midway, above, font=\small] {Score};
\end{tikzpicture}
\caption{RLHF training pipeline.}
\label{fig:rlhf_pipeline}
\end{figure}

\subsubsection{Stage 1: Supervised Fine-Tuning (SFT)}

Fine-tune base model on high-quality demonstrations:
\begin{equation}
\mathcal{L}_{SFT} = -\sum_i \log p(y_i | x_i; \theta)
\end{equation}

\subsubsection{Stage 2: Reward Model Training}

Train model to predict human preferences:
\begin{equation}
\mathcal{L}_{RM} = -\E_{(x, y_w, y_l) \sim \mathcal{D}}\left[\log \sigma(r(x, y_w) - r(x, y_l))\right]
\end{equation}

where $y_w$ is preferred over $y_l$.

\subsubsection{Stage 3: RL Fine-Tuning}

Optimize policy using PPO:
\begin{equation}
\mathcal{L}_{PPO} = \E\left[r(x, y) - \beta \cdot \text{KL}(\pi_\theta(y|x) \| \pi_{ref}(y|x))\right]
\end{equation}

KL penalty prevents diverging too far from SFT model.

\subsection{Direct Preference Optimization (DPO)}

\textbf{Key insight}: Can optimize directly on preferences without explicit reward model.

\begin{equation}
\mathcal{L}_{DPO} = -\E\left[\log \sigma\left(\beta \log \frac{\pi_\theta(y_w|x)}{\pi_{ref}(y_w|x)} - \beta \log \frac{\pi_\theta(y_l|x)}{\pi_{ref}(y_l|x)}\right)\right]
\end{equation}

\textbf{Advantages over RLHF}:
\begin{itemize}
    \item No reward model needed
    \item No RL training instability
    \item Simpler implementation
    \item Often similar results
\end{itemize}

\begin{warningbox}
Tokenization mismatches between training and inference are a silent killer. If you fine-tune a model with one tokenizer but serve with another (or a different version), the model will produce garbage. Always version-pin your tokenizer with your model checkpoint.
\end{warningbox}

\begin{warningbox}
Don't assume bigger models are always better for your task. A fine-tuned BERT-base (110M params) often outperforms zero-shot GPT-4 on domain-specific classification tasks with even modest labeled data (1K+ examples). The right architecture depends on your data budget, latency requirements, and task type.
\end{warningbox}

\begin{interviewq}
\textbf{Q: BERT vs GPT for a new NLP task --- how do you decide?}

\textbf{What the interviewer is testing:} Understanding of encoder vs decoder architectures and practical decision-making.

\textbf{A:} The key question is: \textit{is the task about understanding or generation?}

\begin{itemize}
    \item \textbf{Use encoder (BERT/RoBERTa/DeBERTa) when}: Classification, NER, extractive QA, semantic similarity. These tasks need bidirectional context and produce fixed-size outputs. Encoders are also 5-10$\times$ faster at inference for these tasks.

    \item \textbf{Use decoder (GPT/LLaMA) when}: Text generation, summarization, dialogue, translation, or when you want zero-shot/few-shot capability. Also when you have very little or no labeled data (in-context learning).

    \item \textbf{Use encoder-decoder (T5/BART) when}: Sequence-to-sequence tasks like translation, abstractive summarization, or when you want to frame everything as text-to-text.
\end{itemize}

\textbf{Modern nuance}: With instruction-tuned LLMs, the lines blur. GPT-4 can do classification well zero-shot, but a fine-tuned DeBERTa-base will be 100$\times$ cheaper and often more accurate.

\textbf{Common follow-ups:} ``Can you use a decoder for classification?'' (Yes, but it's inefficient --- you're autoregressing to produce ``positive'' or ``negative'' when you could just do a forward pass.) ``When would you use T5 over GPT?'' (When you need the encoder's bidirectional understanding for the input, like in translation.)
\end{interviewq}

\begin{interviewq}
\textbf{Q: How does in-context learning work? Why can GPT-style models learn from examples in the prompt without gradient updates?}

\textbf{What the interviewer is testing:} Understanding of emergent capabilities and the mechanics of prompt-based learning.

\textbf{A:} In-context learning is still not fully understood, but the leading explanations are:

\begin{enumerate}
    \item \textbf{Implicit Bayesian inference}: The model has seen many task demonstrations during pre-training. Given examples in the prompt, it infers the task distribution and applies the appropriate ``skill'' already encoded in its weights.

    \item \textbf{Attention as gradient descent}: Recent work shows that transformer attention layers can implement a form of gradient descent on the in-context examples. The key-value pairs act as a form of ``training data'' that attention uses to update the representation.

    \item \textbf{Task location, not task learning}: The model has already learned many tasks during pre-training. The examples in the prompt help the model ``locate'' the right task in its weight space, rather than learning something new.
\end{enumerate}

\textbf{Practical implications}: More examples generally help (up to a point), example order matters (recency bias), and the format of examples matters more than you'd expect.

\textbf{Common follow-ups:} ``Why does example order matter?'' (Recency bias in autoregressive models --- later examples have more influence.) ``How many examples do you need?'' (Typically 3-8; beyond that, diminishing returns.)
\end{interviewq}